% Лекц 12. Хотгор, гүдгэр функц ба Тайлорын теорем
% Эх: latex/lectures/lec12-concavity-taylor.tex (scripts/latex_to_beamer.py-ээр хөрвүүлсэн)
% Frame бүр слайдад нэг (эсвэл allowframebreaks-аар хэд хэдэн) хуудас болно;
% хэвлэх хувилбарт frame-ийн хил харагдахгүй, текст үргэлжилсэн байна.

\MathLecture{12}{Хотгор, гүдгэр функц ба Тайлорын теорем}

\section{Хотгор ба гүдгэр функц}
\begin{frame}[allowframebreaks]{Хотгор ба гүдгэр функц}
\begin{definition}
Функц $y=f(x)$ нь $I=(a;b)$ завсар дээр 
$ \forall\,x_1,x_2\in I,\; \lambda\in[0,1]$ хувьд 
\[
  f(\lambda x_1 + (1-\lambda)x_2)
  \geq \lambda f(x_1) + (1-\lambda)f(x_2) \] тэнцэтгэл бишийг хангаж байвал $y=f(x)$ функцийг $I$ завсар дээр 
  \textbf{гүдгэр} (\emph{convex}) функц гэнэ. Хэрэв 
  \[
  f(\lambda x_1 + (1-\lambda)x_2)
  \leq \lambda f(x_1) + (1-\lambda)f(x_2), 
  \] тэнцэтгэл бишийг хангаж байвал $y=f(x)$ функцийг $I$ завсар дээр 
  \textbf{хотгор} (\emph{concave}) функц гэнэ.
\end{definition}
\end{frame}

\begin{frame}[allowframebreaks]{Хотгор ба гүдгэр функц}
\begin{example}
$f(x)=x^2$ функц дурын $I$ завсар дээр хотгор байхыг харуулая. Тодорхойлолтоор 
\[
f(\lambda x_1+(1-\lambda)x_2)
\leq
\lambda f(x_1)+(1-\lambda)f(x_2),
\qquad
x_1,x_2\in I,\ \lambda\in[0,1].
\]
байх ёстой.  $f(x)=x^2$ тул
\[
(\lambda x_1+(1-\lambda)x_2)^2
\le
\lambda x_1^2+(1-\lambda)x_2^2.
\]
Илэрхийлэлийг задлабал
\[
\lambda^2x_1^2+2\lambda(1-\lambda)x_1x_2+(1-\lambda)^2x_2^2
\le
\lambda x_1^2+(1-\lambda)x_2^2.
\]

Баруун талыг зүүн талаас хасвал:
\[
0
\le
\lambda(1-\lambda)(x_1-x_2)^2.
\]

Энэ нь бүх $x_1,x_2$ ба $\lambda\in[0,1]$ үед илт эерэг буюу $0$-оос их тэнцүү байна.  
Иймд тодорхойлолтоор
\[
\boxed{f(x)=x^2 \text{ нь бүх } \mathbb{R}\text{-д хотгор функц юм.}}
\]
\end{example}
\begin{center}
\begin{fitbox}\begin{tikzpicture}[scale=1.0]
  % Coordinate axes
  \draw[->] (-2.5,0) -- (2.5,0) node[right] {$x$};
  \draw[->] (0,-0.5) -- (0,3) node[above] {$y$};

  % Function y = x^2
  \draw[domain=-2:2.2, smooth, variable=\x, thick, blue]
    plot ({\x},{\x*\x}) node[right] {$y = x^2$};

  % Points x1 and x2
  \coordinate (A) at (-0.5,0.25);
  \coordinate (B) at (2,4);
  \fill (A) circle (1.5pt) node[left] {$A(x_1, f(x_1))$};
  \fill (B) circle (1.5pt) node[above right] {$B(x_2, f(x_2))$};

  % Chord between A and B
  \draw[very thick, red] (A) -- (B) node[midway, above, sloped] {хөвч};

  % Convex point
  \coordinate (C) at (1,1);
  \fill (C) circle (1.5pt) node[below right] {$C(\lambda x_1+(1-\lambda)x_2, f(\lambda x_1+(1-\lambda)x_2))$};
  \coordinate (D) at (1,2.5);
  % Dashed line showing function below chord
  \fill (D) circle (1.5pt) node[right] {$D(\lambda x_1+(1-\lambda)x_2,\lambda f(x_1)+(1-\lambda)f(x_2))$};
  \draw[dashed,red] (1,0) -- (D) node[midway, right] {};
\end{tikzpicture}\end{fitbox}
\end{center}
\end{frame}

\begin{frame}[allowframebreaks]{Хотгор ба гүдгэр функц}
\begin{theorem}[Хотгор, гүдгэр ба өсөлт бууралт]
    Хэрэв дифферинциалчлагдах $y=f(x)$ функц $(a;b)$  завсар дээр хотгор (гүдгэр ) бол түүний  уламжлал $(a;b)$ завсар дээр өснө (буурна).
\end{theorem}
\end{frame}

\begin{frame}[allowframebreaks]{Хотгор ба гүдгэр функц}
\begin{theorem}[Хотгор, гүдгэрийг тодорхойлоход уламжлал хэрэглэх]
Хэрэв $(a;b)$  завсар дээр $f''(x)$ хоёр дахь уламжлал оршин байдаг ба $\forall x\in (a;b)$ хувьд
\[
\begin{cases}
f''(x)>0 &\Rightarrow f \text{ хотгор},\\[4pt]
f''(x)<0 &\Rightarrow f \text{ гүдгэр}.
\end{cases}
\]
\end{theorem}

\begin{prof}
Лагранжийн дундаж утгын теорем ашиглая.
$x_1<x_2<x_3$ гурван цэг авбал,
\[
f'(x_2)-f'(x_1)=f''(c_1)(x_2-x_1),\qquad
f'(x_3)-f'(x_2)=f''(c_2)(x_3-x_2).
\]
Хэрэв $f''>0$ бол $f'$ өснө, иймээс графикийн шүргэгчүүдийн налалт өсдөг буюу график гүдгэр байна. $f''<0$ үед эсрэгээр.
\end{prof}
\end{frame}

\begin{frame}[allowframebreaks]{Хотгор ба гүдгэр функц}
\begin{definition}
    $y=f(x)$ функц $(a;c)$ завсар дээр хотгор (гүдгэр ) , $(c;b)$ завсар дээр  гүдгэр  (хотгор)  бол $x=c$ цэгийг түүний нугаралтын цэг гэнэ.
\end{definition}
\end{frame}

\begin{frame}[allowframebreaks]{Хотгор ба гүдгэр функц}
\begin{example}
$f(x)=x^3$ функцийн $f''(x)=6x$ тул $(-\infty,0)$ д гүдгэр, $(0,+\infty)$ д хотгор байна. Иймд $x=0$ нугаралтын цэг байна.
\end{example}

\begin{center}
\begin{fitbox}\begin{tikzpicture}[scale=0.9]
  \draw[->] (-2,0)--(2,0) node[right]{$x$};
  \draw[->] (0,-3)--(0,3) node[above]{$y$};
  \draw[domain=-1.3:1.3,smooth,variable=\x,blue,thick]
    plot({\x},{\x*\x*\x});
  \node at (1.1,1.2){хотгор};
  \node at (-1.1,-1.2){гүдгэр};
  \coordinate (C) at (0,0);
  \fill (C) circle (1.5pt) node[below right]{нугаралт};
\end{tikzpicture}\end{fitbox}
\end{center}
\end{frame}

\section{локал экстремум ба хотгор/гүдгэр}
\begin{frame}[allowframebreaks]{локал экстремум ба хотгор/гүдгэр}
Локал максимумын цэгийн орчимд функц гүдгэр байх нь илэрхий. Иймд сэжигтэй цэгийн орчимд гүдгэр хотгор байгаагаар нь түүний ямар төрлийн цэг болохыг тодорхойлож болно.  
$x_0$ цэг нь $y=f(x)$ функцийн сэжигтэй цэг байг:
\begin{itemize}
  \item Хэрэв $f''(x_0)>0$ бол $x_0$ — локал \textbf{минимум}.
  \item Хэрэв $f''(x_0)<0$ бол $x_0$ — локал \textbf{максимум}.
  \item Хэрэв $f''(x_0)=0$, $f'''(x_0)\ne0$ бол $x_0$ — \textbf{нугаралтын цэг}.
\end{itemize}
\end{frame}

\begin{frame}[allowframebreaks]{локал экстремум ба хотгор/гүдгэр}
\begin{example}
$f(x)=x^4-4x^2$.
\[
f'(x)=4x^3-8x,\qquad f''(x)=12x^2-8.
\]
$f'(x)=0\Rightarrow x=0,\pm\sqrt2$.  
$f''(0)=-8<0$ → максимум,  
$f''(\pm\sqrt2)=16>0$ → минимум.
\end{example}

\begin{center}
\begin{fitbox}\begin{tikzpicture}[scale=0.8]
  \draw[->] (-3,0)--(3,0) node[right]{$x$};
  \draw[->] (0,-5)--(0,5) node[above]{$y$};
  \draw[domain=-2.2:2.2,smooth,variable=\x,blue,thick]
    plot({\x},{\x*\x*\x*\x - 4*\x*\x});
  \node at (-2,-4.2){мин};
  \node at (0,0.5){макс};
  \node at (2,-4.2){мин};
\end{tikzpicture}\end{fitbox}
\end{center}
\end{frame}

\subsection*{3. Тэйлорын теорем}
\begin{frame}[allowframebreaks]{3. Тэйлорын теорем}
\begin{theorem}[Тэйлор]
Хэрэв $f(x)$ нь $a$ цэгийн орчинд $(n+1)$ удаа тасралтгүй уламжлалтай бол $a$ цэгийн орчний $x$ цэгийн хувьд
\begin{equation}\label{Taylor ex}
    f(x)=f(a)+f'(a)(x-a)+\frac{f''(a)}{2!}(x-a)^2+\cdots+\frac{f^{(n)}(a)}{n!}(x-a)^n+R_n(x),
\end{equation}
энд 
\[
R_n(x)=\frac{f^{(n+1)}(c)}{(n+1)!}(x-a)^{n+1},\quad c\in(a,x)
\] байна. 

Хэрэв $x-a=h$ гэж тэмдэглэвэл \[
f(a+h)=f(a)+f'(a)h+\frac{f''(a)}{2!}h^2+\cdots+\frac{f^{(n)}(a)}{n!}h^n+R_n(x),
\] болно. $R_n(x)$-г үлдэгдэл гишүүний Лагранжийн хэлбэр гэнэ.

Энэхүү \ref{Taylor ex} задаргааг $a$ төвтэй Тэйлорын задаргаа гэдэг. 
\end{theorem}

\begin{prof}
$f(x)-P_n(x)$ ба $(x-a)^{n+1}$ хоёр функц $x=a$ дээр 0 утгатай.  
Лагранжийн дундаж утгын теоремийг $(n+1)$ удаа хэрэглэснээр
$a<c<x$ тийм $c$ оршин:
\[
\frac{f(x)-P_n(x)}{(x-a)^{n+1}}=\frac{f^{(n+1)}(c)}{(n+1)!}.
\]
\end{prof}
\end{frame}

\begin{frame}[allowframebreaks]{3. Тэйлорын теорем}
\begin{example}
$f(x)=\sqrt{1+x}$, $a=0$ гэж авъя.  
$f'(x)=\dfrac1{2\sqrt{1+x}},\;
f''(x)=-\dfrac1{4(1+x)^{3/2}}.$  
Тэгвэл
\[
f(x)\approx1+\frac{x}{2}-\frac{x^2}{8}.
\]
$\sqrt{1.1}\approx1+\frac{0.1}{2}-\frac{0.01}{8}=1.04875.$  
Бодит утга $1.0488$ — сайн дөхөлт.
\end{example}
\end{frame}

\begin{frame}[allowframebreaks]{3. Тэйлорын теорем}
\begin{remark}[Маклорены задаргаа]
\ref{Taylor ex} томьёонд $a=0$ үед гарах задаргааг Маклорены задаргаа гэдэг. 
\[
f(x)=\sum_{k=0}^{n}\frac{f^{(k)}(0)}{k!}\,x^k+ \frac{f^{(n+1)}(c)}{(n+1)!}\,x^{n+1}.
\]
Хэрэв $n\to\infty$ үед үлдэгдэл $R_{n}(x)\to 0$ бол Маклорены \emph{цуваа} гэдэг.
\end{remark}
\end{frame}

\begin{frame}[allowframebreaks]{3. Тэйлорын теорем}
\begin{example}[Алдааны үнэлгээ (Лагранж үлдэгдэл)]
$f(x)=e^x$, $a=0$, $n=3$ гэж авъя. Тэгвэл
\[
e^x=1+x+\frac{x^2}{2!}+\frac{x^3}{3!}+\underbrace{\frac{e^c}{4!}x^4}_{R_4(x)},\qquad c\in(0,x).
\]
Жишээ нь $x=0.1$ үед $|R_4(0.1)|=\dfrac{e^{c}}{24}\,0.1^4\le \dfrac{e^{0.1}}{24}\,0.1^4<2.1\times 10^{-6}=0.0000021$.
\end{example}
\end{frame}

\section{Үндсэн функцүүдийн Маклорены задаргаа}
\begin{frame}[allowframebreaks]{Үндсэн функцүүдийн Маклорены задаргаа}
Дараах задаргааг \ref{Taylor ex} ашиглан хялбар шалгаж болно.
\[
e^x \;=\; 1 + x + \frac{x^2}{2!} + \cdots + \frac{x^n}{n!} \;+\; R_{n+1}(x), 
\qquad x\in\mathbb{R}\]
\[R_{n+1}(x)=\frac{e^{c}}{(n+1)!}\,x^{\,n+1},\ \ c\in(0,x).
\]
\[
\sin x \;=\; x - \frac{x^3}{3!} + \frac{x^5}{5!} - \cdots + 
(-1)^{m}\frac{x^{2m+1}}{(2m+1)!} \;+\; R_{2m+1}(x), \qquad x\in\mathbb{R}
\]
\[
R_{2m+1}(x)=\frac{\sin^{(2m+2)}(c)}{(2m+2)!}\,x^{\,2m+2},\ \ c\in(0,x).
\]
\[
\cos x \;=\; 1 - \frac{x^2}{2!} + \frac{x^4}{4!} - \cdots + 
(-1)^{m}\frac{x^{2m}}{(2m)!} \;+\; R_{2m}(x), \qquad x\in\mathbb{R}
\]
\[
R_{2m}(x)=\frac{\cos^{(2m+1)}(c)}{(2m+1)!}\,x^{\,2m+1},\ \ c\in(0,x).
\]
\[
\ln(1+x) \;=\; x - \frac{x^2}{2} + \frac{x^3}{3} - \cdots + 
(-1)^{n-1}\frac{x^{n}}{n} \;+\; R_{n}(x), \qquad x\in(-1;1)
\]
\[
R_{n+1}(x)=\frac{\ln^{(n+1)}(1+c)}{(n+1)!}\,x^{\,n+1}
=\frac{(-1)^{n}\,n!}{(n+1)!}\,\frac{x^{\,n+1}}{(1+c)^{n+1}}
= (-1)^{n}\,\frac{x^{\,n+1}}{(n+1)\,(1+c)^{n+1}},\ \ c\in(0,x).
\]
\[
\arctan x \;=\; x - \frac{x^3}{3} + \frac{x^5}{5} - \cdots + 
(-1)^{m}\frac{x^{2m+1}}{2m+1} \;+\; R_{2m+1}(x), \qquad x\in(-1;1)
\]
\[
R_{2m+1}(x)=\frac{(\arctan)^{(2m+2)}(c)}{(2m+1)!}\,x^{\,2m+1},\ \ c\in(0,x).
\]
\[
(1+x)^{\alpha} \;=\; 1 + \alpha x + \frac{\alpha(\alpha-1)}{2!}x^2 + \cdots +
\frac{\alpha(\alpha-1)\cdots(\alpha-n+1)}{n!}\,x^n \;+\; R_{n}(x), \qquad x\in(-1;1)
\]
\[
R_{n}(x)=\frac{(1+x)^{\alpha\,(n+1)}(c)}{(n+1)!}\,x^{\,n+1}
=\frac{\alpha(\alpha-1)\cdots(\alpha-n)}{(n+1)!}\,(1+c)^{\alpha-n-1}\,x^{\,n+1},\ \ c\in(0,x).
\]
Дээрх задаргаанд $\alpha=-1, x=-t, n=\infty$ гэвэл 
\[
\frac{1}{1-t} \;=\; 1 + t + t^2 + \cdots + t^n +\cdots 
\]
болж бидний мэдэх төгсгөлгүй буурах геометрийн прогрессийн нийлбэр гарна. Бусад задаргаанд ч мөн адил $n=\infty$ авч үзвэл цуваа үүснэ.
\end{frame}

\begin{frame}[allowframebreaks]{Үндсэн функцүүдийн Маклорены задаргаа}
\begin{example}[Тэйлор $a=\pi/4$ төвтэй]
$f(x)=\sin x$ функцийн $a=\pi/4$ төвтэй $n=2$ задаргааг бич.

$f(\tfrac{\pi}{4})=\tfrac{\sqrt2}{2}$,\ $f'(\tfrac{\pi}{4})=\tfrac{\sqrt2}{2}$,\ $f''(\tfrac{\pi}{4})=-\tfrac{\sqrt2}{2}$ тул
\[
\sin x=\frac{\sqrt2}{2}+\frac{\sqrt2}{2}(x-\tfrac{\pi}{4})-\frac{\sqrt2}{4}(x-\tfrac{\pi}{4})^2+\frac{-\sin(c)}{3!}(x-\tfrac{\pi}{4})^3.
\] Хэрэв энэ задаргааг ашиглах болбол ихэвчлэн үлдэгдэл гишүүнийг орхин ойролцоо тооцоолол хийдэг. Учир нь $c$ тоо олдох л болохоос хэрхэн олохыг мэдэхгүй.
\end{example}
\end{frame}

\begin{frame}[allowframebreaks]{Үндсэн функцүүдийн Маклорены задаргаа}
\begin{example}[Маклорен 3-р зэргийн ойролцоо] $y=\ln(1+x)$ функцийг 3 зэргийн олон гишүүнтээр ойролцоогоор соль.\\
$\ln(1+x)=x-\dfrac{x^2}{2}+\dfrac{x^3}{3}+R_4(x)$,\quad 
$|R_4(x)|=|- \dfrac{x^4}{4(1+c)^4}|\leq \dfrac{x^4}{4}$ тул  $|\ln(1+x) - (x-\tfrac{x^2}{2}+\tfrac{x^3}{3})|\lesssim \dfrac{|x|^4}{4}$.
\end{example}
\end{frame}

\begin{frame}[allowframebreaks]{Үндсэн функцүүдийн Маклорены задаргаа}
\begin{example}[Тэйлороор ойролцоо бодох]
$\sqrt{1+x}=(1+x)^{1/2}=1+\frac{x}{2}-\frac{x^2}{8}+\frac{x^3}{16}+\cdots$ ($|x|<1$). 
Жишээ нь $\sqrt{9.1}=\sqrt{9}\sqrt{1+0.1/9}\approx 3\!\left(1+\frac{0.1}{18}-\frac{(0.1)^2}{8\cdot 9^2}\right)$.
\end{example}
\end{frame}

\section{Функц шинжлэн график байгуулах 4 алхам}
\begin{frame}[allowframebreaks]{Функц шинжлэн график байгуулах 4 алхам}
Уламжлал ашиглан функцийн графикийг байгуулахад нийлээд нарийн зурах боломжтой. Дараах алхамуудын дагуу гарфик байгуулдаг.
\end{frame}

\subsection*{Алхам 1. $f(x)$-ийн ерөнхий шинжилгээ (графикийн үндсэн мэдээлэл)}
\begin{frame}[allowframebreaks]{Алхам 1. $f(x)$-ийн ерөнхий шинжилгээ (графикийн үндсэн мэдээлэл)}
\begin{itemize}
  \item \textbf{Тодорхойлогдох муж} $D_f$: томъёоны хязгаарлалт (хориглогдох тоон утга, тэнцэтгэл биш).
  \item \textbf{Тэгш/Сондгой \& Үе}: 
  \[
    f(-x)\stackrel{?}{=}f(x)\;\;(\text{тэгш}), \qquad
    f(-x)\stackrel{?}{=}-f(x)\;\;(\text{сондгой}), \qquad
    f(x+T)\stackrel{?}{=}f(x).
  \]
  \item \textbf{Тэнхлэгтэй огтлол \& тэмдэг}: $f(x)=0$ тэгүүд, $\;f(x)\gtrless 0$ тэмдэгийн муж (тэнхлэгийн дээгүүр/доогуур).
  \item \textbf{Хязгаарууд (шаардлагатай бол)}: $x\to a,\;\pm\infty$ үед $\lim f(x)$, хэвтээ/босоо/ ташуу асимптот боломж.
\end{itemize}
\end{frame}

\subsection*{Алхам 2. Нэгдүгээр уламжлал: өсөлт/бууралт, локал экстеомум}
\begin{frame}[allowframebreaks]{Алхам 2. Нэгдүгээр уламжлал: өсөлт/бууралт, локал экстеомум}
\begin{itemize}
  \item \textbf{Сэжигтэй цэгүүд}: $f'(x)=0$ болон $f'$ оршихгүй цэгүүд.
  \item \textbf{Өсөлт/Бууралт}: $f'(x)\geq 0$ тэмдэгийн хүснэгт байгуулж завсар бүрт монотон чанарыг тогтооно.
  \item \textbf{Локал экстремум}: $f'$-ийн тэмдэг солигдох байдлаар local \(\max/\min\) цэгүүдийг олно.
\end{itemize}
\end{frame}

\subsection*{Алхам 3. Хоёрдугаар уламжлал: хотгор/гүдгэр, нугаралтын цэг}
\begin{frame}[allowframebreaks]{Алхам 3. Хоёрдугаар уламжлал: хотгор/гүдгэр, нугаралтын цэг}
\begin{itemize}
  \item \textbf{Хотгор/Гүдгэр}: $f''(x)\gtrless 0$ тэмдэгийн хүснэгт байгуулж завсар бүрийн хэлбэрийг тогтооно.
  \item \textbf{Нугаралтын цэг}: $f''(x)=0$ буюу $f''$ оршихгүй ба  хотгор гүдгэр солигдох цэгүүд.
\end{itemize}
\end{frame}

\subsection*{Алхам 4. Нэгтгэх хүснэгт \& График}
\begin{frame}[allowframebreaks]{Алхам 4. Нэгтгэх хүснэгт \& График}
\begin{itemize}
  \item Алхам 1–3 дахь бүх мэдээллээ нэг хүснэгтэд нэгтгээд,
  \item Тэгш/сондгой байдал, асимптот, $x$-тэнхлэгийг огтлох цэгүүд, өсөлт/бууралт, хотгор/гүдгэр,
        экстремум, нугаралтын цэгүүдийг тэмдэглэн \textbf{график} зурах.
\end{itemize}
\end{frame}

\subsection*{Нэгтгэх хүснэгтийн жишиг загвар}
\begin{frame}[allowframebreaks]{Нэгтгэх хүснэгтийн жишиг загвар}
\begin{center}
\renewcommand{\arraystretch}{1.2}
\begin{fitbox}\begin{tabular}{|c|c|c|c|c|c|}
\hline
Завсар & $f(x)$-ийн тэмдэг & $f'(x)$ тэмдэг & Монотон & $f''(x)$ тэмдэг & Хотгор/Гүдгэр \\
\hline
$(\;\;,\;\;)$ & $+$/$-$ & $+$/$-$ & Өснө/Буурна & $+$/$-$ & Гүдгэр/Хотгор \\
\hline
$\cdots$ & $\cdots$ & $\cdots$ & $\cdots$ & $\cdots$ & $\cdots$ \\
\hline
\end{tabular}\end{fitbox}
\end{center}
\end{frame}

\subsection*{Алхам 1. Ерөнхий шинжилгээ}
\begin{frame}[allowframebreaks]{Алхам 1. Ерөнхий шинжилгээ}
\begin{itemize}
  \item \textbf{Тодорхойлогдох муж:} $D_f=\mathbb{R}$.
  \item \textbf{Тэнхлэгтэй огтлол:} $f(0)=(0+1)e^0=1$ → $y$-тэнхлэгийг $(0,1)$ цэгт огтлоно.
  \item $f(x)=0 \ \Rightarrow\ (x+1)e^x=0 \ \Rightarrow\ x=-1$ → $x=-1$ дээр $x$-тэнхлэгийг огтлоно.
  \item \textbf{Тэмдэг:} $f(x)<0$ ⇔ $x<-1$, $f(x)>0$ ⇔ $x>-1$.
  \item \textbf{Хязгаарууд:} 
  \[
  \lim_{x\to -\infty} (x+1)e^x = 0^{-}, \qquad 
  \lim_{x\to +\infty} (x+1)e^x = +\infty.
  \]
\end{itemize}
\end{frame}

\subsection*{Алхам 2. Нэгдүгээр уламжлал (өсөлт, бууралт, экстремум)}
\begin{frame}[allowframebreaks]{Алхам 2. Нэгдүгээр уламжлал (өсөлт, бууралт, экстремум)}
\[
f'(x) = (x+1)'e^x + (x+1)(e^x)' = e^x + (x+1)e^x = e^x(x+2).
\]
\begin{itemize}
  \item $f'(x)=0$ ⇔ $x=-2$.
  \item $e^x>0$ бүх $x$-д, тэгэхээр $f'(x)$-ийн тэмдэг нь $(x+2)$-оор тодорхойлогдоно.
  \[
  f'(x)<0 \text{  үед } x<-2, \qquad f'(x)>0 \text{  үед } x>-2.
  \]
  \item $\Rightarrow$ $x=-2$ нь \textbf{local минимум}.
  \item $f(-2)=(-1)e^{-2}=-\dfrac{1}{e^2}\approx -0.135$.
\end{itemize}
\end{frame}

\subsection*{Алхам 3. Хоёрдугаар уламжлал (хотгор, гүдгэр, нугаралт)}
\begin{frame}[allowframebreaks]{Алхам 3. Хоёрдугаар уламжлал (хотгор, гүдгэр, нугаралт)}
\[
f''(x) = \frac{d}{dx}[e^x(x+2)] = e^x(x+2)+e^x=e^x(x+3).
\]
\begin{itemize}
  \item $f''(x)=0$ ⇔ $x=-3$ → нугаралтын цэг.
  \item $f''(x)<0$ ($x<-3$) → хотгор, $f''(x)>0$ ($x>-3$) → гүдгэр.
  \item $\Rightarrow$ $x=-3$ нь нугаралтын цэг.
\end{itemize}
\end{frame}

\subsection*{Алхам 4. Мэдээллийг нэгтгэн хүснэгтэд оруулж график байгуулах}
\begin{frame}[allowframebreaks]{Алхам 4. Мэдээллийг нэгтгэн хүснэгтэд оруулж график байгуулах}
\begin{center}
\renewcommand{\arraystretch}{1.3}
\begin{fitbox}\begin{tabular}{|c|c|c|c|c|c|}
\hline
Завсар & $f(x)$ тэмдэг & $f'(x)$ тэмдэг & Өсөлт/Бууралт & $f''(x)$ тэмдэг & Хотгор/Гүдгэр \\
\hline
$(-\infty,-3)$ & $-$ & $-$ & Буурна & $-$ & Хотгор \\
\hline
$(-3,-2)$ & $-$ & $-$ & Буурна & $+$ & Гүдгэр \\
\hline
$(-2,\infty)$ & $+$ & $+$ & Өснө & $+$ & Гүдгэр \\
\hline
\end{tabular}\end{fitbox}
\end{center}
\end{frame}

\subsection*{График}
\begin{frame}[allowframebreaks]{График}
\begin{center}
\begin{fitbox}\begin{tikzpicture}[scale=1.0]
  % Тэнхлэгүүд
  \draw[->] (-5,0)--(3,0) node[below right] {$x$};
  \draw[->] (0,-2)--(0,5) node[above left] {$y$};

  % Функцийн муруй
  \draw[smooth,domain=-5:1,samples=200,thick,blue] plot(\x,{(\x+1)*exp(\x)});

  % Тэмдэглэсэн цэгүүд
  \fill (-2,{-1/exp(2)}) circle (2pt) node[below left] {$(-2,-\frac{1}{e^2})$};
  \fill (0,1) circle (2pt) node[above right] {$(0,1)$};
  \fill (-1,0) circle (2pt) node[below left] {$( -1,0)$};
  \draw[dashed] (-2,0)--(-2,{-1/exp(2)})--(0,{-1/exp(2)});

  % Тайлбарууд
  \node at (-4,3.5) {\small $f(x)=(x+1)e^x$};
  \node[blue] at (-2.5,-1.4) {\small Мин цэг};
  \node[blue] at (-3.2,0.8) {\small Хотгор};
  \node[blue] at (-1,3) {\small Гүдгэр};
\end{tikzpicture}\end{fitbox}
\end{center}
\end{frame}

% Семинар ба гэрийн даалгавар (Берманы бодлогын хураамжаас)
% Лекц 12: семинар ба гэрийн даалгавар. Бодлогууд: Г.Н. Берман, "Сборник задач по курсу
% математического анализа" (22-р хэвлэл). [Б.N] нь номын бодлогын дугаар.
% Тэмдэглэл: Берманы "вогнутый" = манай "гүдгэр" ($f''>0$), "выпуклый" = "хотгор" ($f''<0$).

\ProblemSet{Семинар 12}
\begin{frame}[allowframebreaks]{Семинар 12}
\pgroup{Хоёрдугаар эрэмбийн уламжлал ашиглан экстремумыг ол.}
\begin{plist}
\prob{1267}{$y=x^3-2ax^2+a^2x$ \ ($a>0$)}
\prob{1273}{$y=x^2e^{-x}$}
\end{plist}
\pgroup{Гүдгэр, хотгор завсар ба нугаралтын цэг.}
\begin{plist}
\prob{1281}{$4y=x\arctg x$ функцийн график хаа сайгүй гүдгэр болохыг харуул.}
\prob{1282}{$y=\ln(x^2-1)$ функцийн график хаа сайгүй хотгор болохыг харуул.}
\prob{1287}{$y=x^3-5x^2+3x-5$}
\prob{1289}{$y=x^4-12x^3+48x^2-50$}
\prob{1296}{$y=\ln(1+x^2)$}
\prob{1299}{$y=e^{\arctg x}$}
\prob{1305}{$a$ ба $b$ ямар утгатай үед $(1,3)$ цэг $y=ax^3+bx^2$ муруйн нугаралтын цэг болох вэ?}
\end{plist}
\pgroup{Асимптотыг ол.}
\begin{plist}
\prob{1377}{$y=\frac{1}{x^2-4x+5}$}
\prob{1379}{$2y(x+1)^2=x^3$}
\prob{1387}{$y=xe^x$}
\end{plist}
\pgroup{Функцийг бүрэн судалж графикийг байгуул.}
\begin{plist}
\prob{1398}{$y=\frac{x}{1+x^2}$}
\prob{1400}{$y=\frac{x}{x^2-1}$}
\prob{1408}{$y=\frac{x^3}{3-x^2}$}
\prob{1419}{$y=x-\ln(x+1)$}
\end{plist}
\pgroup{Тейлорын томьёо.}
\begin{plist}
\prob{1498}{$x^4-5x^3+x^2-3x+4$ олон гишүүнтийг $x-4$-ийн зэргээр задал.}
\prob{1504}{$y=xe^x$ функцийн $x_0=0$ цэг дээрх $n$-р эрэмбийн Тейлорын (Маклорены) томьёог бич.}
\prob{1505}{$y=\sqrt x$ функцийн $x_0=4$ цэг дээрх $n$-р эрэмбийн Тейлорын томьёог бич.}
\prob{1510}{$y=\tg x$ функцийн $x_0=0$ цэг дээрх 2-р эрэмбийн Тейлорын томьёог бич.}
\end{plist}
\end{frame}

\ProblemSet{Гэрийн даалгавар 12}
\begin{frame}[allowframebreaks]{Гэрийн даалгавар 12}
\pgroup{Хоёрдугаар эрэмбийн уламжлал ашиглан экстремумыг ол.}
\begin{plist}
\prob{1268}{$y=x^2(a-x)^2$}
\prob{1270}{$y=x+\sqrt{1-x}$}
\prob{1271}{$y=x\sqrt{2-x^2}$}
\prob{1276}{$a$-ийн ямар утгад $f(x)=a\sin x+\frac13\sin3x$ функц $x=\frac{\pi}{3}$ цэг дээр экстремумтай вэ? Энэ нь максимум уу, минимум уу?}
\prob{1277}{$y=a\ln x+bx^2+x$ функц $x_1=1$, $x_2=2$ цэгүүдэд экстремумтай байх $a,b$-г ол. Энэ үед $x_1$ дээр минимум, $x_2$ дээр максимум болохыг харуул.}
\end{plist}
\pgroup{Нугаралтын цэг, гүдгэр ба хотгор завсрыг ол.}
\begin{plist}
\prob{1288}{$y=(x+1)^4+e^x$}
\prob{1290}{$y=x+36x^2-2x^3-x^4$}
\prob{1292}{$y=(x+2)^6+2x+2$}
\prob{1293}{$y=\frac{x^3}{x^2+3a^2}$ \ ($a>0$)}
\prob{1294}{$y=a-\sqrt[3]{x-b}$}
\prob{1300}{$y=x^4(12\ln x-7)$}
\prob{1301}{$y=\frac{x+1}{x^2+1}$ муруй нэг шулуун дээр орших гурван нугаралтын цэгтэйг харуул.}
\prob{1307}{$a$-ийн ямар утгад $y=e^x+ax^3$ функцийн график нугаралтын цэгтэй вэ?}
\end{plist}
\pgroup{Асимптотыг ол.}
\begin{plist}
\prob{1378}{$y=c+\frac{a^3}{(x-b)^2}$}
\prob{1386}{$y=x\ln\bigl(e+\frac1x\bigr)$}
\prob{1388}{$y=xe^{2/x}+1$}
\prob{1390}{$y=2x+\arctg\frac{x}{2}$}
\end{plist}
\pgroup{Функцийг бүрэн судалж графикийг байгуул.}
\begin{plist}
\prob{1399}{$y=\frac{1}{1-x^2}$}
\prob{1402}{$y=\frac{x^2}{x^2-1}$}
\prob{1405}{$y=\frac1x+4x^2$}
\prob{1409}{$y=\frac{x^3}{2(x+1)^2}$}
\prob{1416}{$y=\frac{x}{e^x}$}
\prob{1420}{$y=\ln(x^2+1)$}
\prob{1423}{$y=xe^{-x^2/2}$}
\prob{1427}{$y=x+\sin x$}
\end{plist}
\pgroup{Тейлорын томьёо.}
\begin{plist}
\prob{1499}{$x^3+3x^2-2x+4$ олон гишүүнтийг $x+1$-ийн зэргээр задал.}
\prob{1503}{$y=\frac1x$ функцийн $x_0=-1$ цэг дээрх $n$-р эрэмбийн Тейлорын томьёог бич.}
\prob{1508}{$y=\sin^2x$ функцийн $x_0=0$ цэг дээрх $2n$-р эрэмбийн Тейлорын томьёог бич.}
\prob{1514}{$y=2x^6-x^3+3$ функцийн $x=0$ цэгийн орчим дахь төлөвийг Тейлорын томьёогоор тогтоо.}
\prob{1528}{$\ln(1+x)\approx x-\frac{x^2}{2}+\frac{x^3}{3}-\frac{x^4}{4}$ томьёогоор $\ln1{,}5$-г олж, алдааг үнэл.}
\end{plist}
\end{frame}

